% SVN info for this file
\svnidlong
{$HeadURL$}
{$LastChangedDate$}
{$LastChangedRevision$}
{$LastChangedBy$}

\chapter*{Note per la lettura}
\labelChapter{introduzione}

\begin{introduction}
‘‘Un matematico è una macchina per trasformare caffè in teoremi.''
\begin{flushright}
	\textsc{Alfréd Rényi,} studioso del teorema di Van Moka-mpen.
\end{flushright}
\end{introduction}

\lettrine[findent=1pt, nindent=0pt]{S}{enza} troppe pretese di formalità, com'è intuibile dal termine dal termine tecnico \textit{manualozzo} e dalle citazioni a inizio capitolo, queste note sono nate come appunti a quattro mani basati sul corso di \textsc{Analisi Matematica 3} tenuto dai docenti Walter Dambrosio e Davide Zucco nell'Anno Accademico 2021-2022 presso il Dipartimento di Matematica dell'Università degli Studi di Torino.\\
I prerequisiti necessari sono gli argomenti trattati nel corso di \textsc{Analisi Matematica Uno}.\\
In aggiunta a ciò, potete trovare a fine libro delle utili \textit{postille} con alcune digressioni interessanti, nonché tabelle ed elenchi riepilogativi dei teoremi, delle definizioni e delle proprietà affrontate.\\
Per quanto ci piacerebbe esserlo, non siamo \textit{esseri infallibili}: ci saranno sicuramente sfuggiti degli errori (o degli \textit{orrori}, la cui causa è solamente degli autori che non hanno studiato bene e non dei professori, chiaramente), per cui vi chiediamo gentilmente di segnalarceli su \textcolor{redill}{\url{https://maxmaci.github.io}} per correggerli e migliorare le future edizioni del \textit{manualozzo}.\\
I disegni sono stati realizzati da Massimo Bertolotti, l'addetto alla grafica e ai capricci di \LaTeX\ (ed è molto capriccioso, fidatevi). Chi volesse dilettarsi può cercare di distinguere chi fra i due autori ha scritto cosa, non dovrebbe essere troppo difficile.
\vfill
\begin{center}
	Prima edizione, compilato il \today.\\
			\includegraphics[trim=0cm 0cm 0cm 0cm,clip,scale=0.5]{images/Cc-by-nc-sa_icon.pdf}\\
	{\footnotesize This work is licensed under a \href{https://creativecommons.org/licenses/by-sa/4.0/}{Attribution-NonCommercial-ShareAlike 4.0 International.}}
\end{center}